\section{The Standing Relation}

Standing is relative. It is never a free-standing property of a candidate and never a universal permission granted by an authoritative evaluator. A statement can be entitled to count as a hypothesis worth testing without being entitled to count as evidence for anything, a model can be entitled to be used for exploration without being entitled to be deployed, a decision can be procedurally authorized while its consequences remain indefensible, and a result can count within the scope in which it was established and cease to count when it is carried beyond that scope. Each of these observations is familiar, and each is usually stated as a qualification added to a verdict. The proposal here is to make the relativity part of the definition, so that the arguments on which a verdict depends are written down wherever the verdict is.

The definitions require a minimal notion of state. A \emph{state} $s$ of a system comprises a shared record of inscriptions with their provenance, a designation of which items are relied upon and for what, and the set $K(s)$ of commitments in force. An \emph{operation} $o$ is a proposed change. A \emph{transition} is a triple $\tau=(s,o,s')$ in which $s'$ is the state that applying $o$ in $s$ would produce. Transitions are state-indexed on purpose: the same operation applied in two different states yields two different transitions, and the second proposition of the next section shows that these can differ in standing even when nothing about the operation or the evidence differs. A \emph{history} is a sequence of transitions in which each begins in the state where the previous one ended, and a transition \emph{occurs} when it is an element of the history the system actually follows. A \emph{rule system} assigns to each state the set of operations it permits, and a transition is \emph{licensed} when its operation is permitted in its initial state.

\begin{definition}[Standing]
Let $\tau=(s,o,s')$ be a transition. Let $C$ be a consequence class, the use or effect the transition would license. Let $S$ be a scope, the set of contexts, domains or populations in which the result would be relied upon. Let $V$ be the evidence available, understood as a body of observations, derivations, records and testimony. Let $P$ be an evaluation procedure, a rule that takes a transition together with $C$, $S$, $V$ and the commitments $K=K(s)$ and returns acceptance or rejection. Let $t$ be the time of assessment, which fixes the version of $V$ and of $P$ that is available. The transition has \emph{standing} when $P$ accepts it, which is written
\[
\Std(\tau \mid C,\,S,\,V,\,P,\,K,\,t).
\]
\end{definition}

The relation is a predicate evaluated relative to a context. It is not a numerical measure, since nothing in the definition assigns degrees, and it is not required to be computable, since $P$ may be a committee, a court or a practice that no algorithm captures. The absence of numbers is deliberate: a number attached to standing would suggest a common scale on which a proof, a clinical trial and a court ruling could be ranked, and the account is built to deny that such a scale exists.

Evidence $V$ and procedure $P$ are separate arguments from the start, and this has a direct consequence. A transition can fail to have standing because the evidence is inadequate for the consequence, because the procedure is defective or ill-suited to the consequence, because it conflicts with a commitment in $K$, or because of any combination of these. A procedure applied to excellent evidence can still produce a verdict that does not deserve to count, and excellent procedure applied to poor evidence cannot repair the poverty of the evidence. Keeping the arguments separate makes it possible to say which has failed, and the two kinds of failure call for different remedies, since better evidence can be gathered while a defective procedure must itself be assessed, which is the subject of the third proposition.

The remaining arguments do their own work. The consequence class $C$ matters because the same candidate may deserve to count for a use with modest effects and not for one with severe or irreversible effects. The scope $S$ matters because a result established in some contexts may not extend to others. The commitments $K$ matter because they change what a transition would newly make irreversible, which is why standing attaches to a transition in a particular state and not to a candidate in the abstract. The time $t$ matters because evidence accumulates, procedures are revised, and what had standing at one time may not at another, without anyone having been wrong at the earlier time.

The relation separates three things that are usually run together. \emph{Occurrence} is what the system did. \emph{Licensing} is what its rules permit. \emph{Standing} is whether, given $V$ and $P$ and the rest, the transition ought to count for the stated purpose. Licensing is a feature of a rule system and standing is a feature of a warrant relation, and in a well-designed system the rules license exactly the transitions that would have standing. No system is well designed in this sense for all transitions in all states, and the gaps run in both directions. Rules can license what has no standing, and can fail to license what would have it. The first proposition below makes this precise, and the later sections return to the gap repeatedly because it is where most failures are located.

There are three senses of standing, distinguished by what plays the roles of $V$, $P$ and $K$. In the \emph{epistemic} sense, the evidence $V$ is a body of observations and arguments, the procedure $P$ is a method of inquiry with norms of inference, replication and error control, and the commitments $K$ are the claims on which a research community has already begun to build. In the \emph{computational} sense, $V$ is a set of execution traces, proofs, test results or formal derivations, $P$ is a checker, verifier or test harness whose behavior is specified in advance, and $K$ is the state of deployed systems and their dependencies. In the \emph{institutional} sense, $V$ is a submitted record, testimony or filing, $P$ is a decision procedure specified by rules and exercised by a body with jurisdiction, and $K$ is the body of prior rulings, ratified decisions and contracts.

The three senses are defined in parallel so that none is made the default. A proof that passes a formal checker has computational standing for the claim that a specified derivation is valid under the checker's rules, and it has epistemic standing for the mathematical claim only to the extent that the formalization corresponds to the claim, which is a separate question. A regulator that accepts a verified computation as part of a filing treats it as evidence $V$ in an institutional procedure, and the regulator's authorization does not become a computational verdict by that route. An argument may pass from one sense to another only through a stated \emph{bridge}, which says what in one sense stands in for $V$, $P$ or $K$ in the other and what is lost in the translation. Statements that move silently across such a boundary, saying that something was verified and therefore authorized, or authorized and therefore true, are the standard form of the error the framework is meant to expose.

Four short contrasts show how much the relativity matters. A claim that a compound reduces a symptom in a pilot study has standing as a hypothesis warranting a larger trial and has not the standing of evidence for prescribing it. A model evaluated on a held-out benchmark has standing for exploratory use in a research setting and not, on that evidence alone, for decisions about individuals. A procurement decision adopted by the proper committee under the proper rules has institutional standing as a decision, and whether its consequences are defensible is a question about a different evidence base and a different procedure. A result established for a laboratory population has standing within the scope of that population and loses it when applied to a population on which the procedure was never run. In each case no one has made a mistake of fact, and the difficulty lies wholly in what the verdict was permitted to carry.

Finally, the relation is something to be assessed and not something that can be looked up. There is no general algorithm that, given a transition and a context, returns whether standing holds, since $P$ may itself be contested and $V$ may be incomplete in ways that cannot be detected from inside the procedure. The relation is useful because it names the arguments that any assessment has to supply, and because it makes visible when an assessment has supplied fewer of them than the verdict requires. The next section states what can be said in general about how the relation behaves.
